

% \begin{itemize}[label=(\\alph*)]
\subsection{Opttimisation} 
    In principle it is possible to determine exactly what x will be at the end of the programme: \\
    it starts at 1, then func adds 1 to it (x = 2) and then main adds 3 to it. Hence when printed it will show the value 5.
    A good optimised code will therefore look like:
    
\begin{lstlisting}
    int main()
    {
        printf("5\n");
    }
\end{lstlisting}

However, when the compiler looks at func.c, it does not know the value of x, so it just produces a function that adds 1 to x. In main, something similar happens; 
as x is a global variable, it cannot guarantee func does not change it (and in fact it does) and thus does not know the value of x after the call to func. Therefore the compiler is forced to code that adds 3
to whatever the value of x is at that moment. In merging the 2 functions into a single file, the compiler can see both modifications, so it does know what the state of x is before the func call, knows the state after the func call and therefore knows the exact value it has to
print at compile time, resulting in the optimised code. 


\subsection{Compiler report}
After running compiling with \texttt{gcc-14 -c -O1 -fopt-info-all *.c}

\includesource{1/gccrep.txt}


